%% toc: P15 162 162.5,162.7,162.8,162.9
\section{Stopping at the first variance exceedance}
\label{sec:annex-p15-w02c-p15-021-first-exceedance}

Throughout, \(X\) is a real square-integrable martingale with
\(X_0=0\) for the filtration \((\mathcal F_n)_{n\geq0}\), and
\(a_n\), \(A_n\), \(A_\infty\) are the versions fixed in
\cref{def:p15-predictable-variation}: each \(a_n\) is finite,
nonnegative and \(\mathcal F_{n-1}\)-measurable, \(A_0=0\) and
\(A_n=a_1+\cdots+a_n\). The level \(c>0\) is fixed, and \(S_c\) and
\(T_c\) are as in \cref{def:p15-w02c-p15-021-first-exceedance}.

\begin{annexitem}{annex:p15-w02c-p15-021-exceedance-shift}{thm:p15-w02c-p15-021-exceedance-shift}{Comparing the two stopping rules}
\begin{proof}
Since \(A_0=0\leq c\), the index \(0\) never belongs to
\(\{n\geq0:A_n>c\}\), so \(S_c\geq1\) and
\[
 S_c=\inf\{m\geq1:A_m>c\}
    =1+\inf\{n\geq0:A_{n+1}>c\}
    =1+T_c,
\]
where the middle equality substitutes \(m=n+1\) and both sides are
infinite together, because the two index sets are empty together.

For \(n\geq0\),
\[
 \{S_c\leq n\}=\bigcup_{j=1}^n\{A_j>c\}.
\]
For \(j\geq1\), \(A_j\) is a finite sum of the
\(\mathcal F_{j-1}\)-measurable variables \(a_1,\ldots,a_j\), so
\(\{A_j>c\}\in\mathcal F_{j-1}\subseteq\mathcal F_n\) when \(j\leq n\).
The union, empty when \(n=0\), therefore lies in \(\mathcal F_n\), and
\(S_c\) is a stopping time.

If \(A_\infty>c\), then, because \(A_\infty=\sup_nA_n\), some \(A_n\)
exceeds \(c\), so \(S_c\leq n<\infty\). If \(S_c=k<\infty\), then
\(A_\infty\geq A_k>c\). Hence \(\{S_c<\infty\}=\{A_\infty>c\}\).

On \(\{S_c<\infty\}\) put \(k=S_c\geq1\). Then \(k-1=T_c\) and
\(A_k=A_{k-1}+a_k=A_{T_c}+a_{S_c}\), and \(A_k>c\) by the definition of
\(S_c\). If \(k=1\), then \(A_{T_c}=A_0=0\leq c\). If \(k\geq2\), then
\(k-1\geq1\) is an index smaller than the infimum \(k\), so
\(A_{k-1}\leq c\). In both cases \(A_{T_c}\leq c\).
\end{proof}
\end{annexitem}

\begin{annexitem}{annex:p15-w02c-p15-021-exceedance-bound}{thm:p15-w02c-p15-021-exceedance-bound}{The idle case and the bounded-increment case}
\begin{proof}
Fix a sample point. Suppose first \(S_c=k<\infty\). For every
\(n\geq k\), \(S_c\wedge n=k\) and \(A_{S_c\wedge n}=A_k>c\). Thus the
bound fails for all \(n\geq k\), and the first statement on the left
of the displayed equivalences is false. Suppose next \(S_c=\infty\).
Then no \(A_n\) exceeds \(c\), so \(A_{S_c\wedge n}=A_n\leq c\) for every
\(n\). This proves the first equivalence. The second is the identity
\(\{S_c<\infty\}=\{A_\infty>c\}\) of
\cref{thm:p15-w02c-p15-021-exceedance-shift}, read through complements.

Since the equivalences hold at every sample point, the event on which
\(A_{S_c\wedge n}\leq c\) for every \(n\) equals \(\{A_\infty\leq c\}\).
The first event has probability one exactly when the second does. On
\(\{S_c=\infty\}\), \(S_c\wedge n=n\), so \(X_{S_c\wedge n}=X_n\) for
every \(n\); this holds almost surely when \(A_\infty\leq c\) almost
surely.

Assume now \(a_n\leq K\) almost surely for every \(n\), with \(K\)
deterministic. The countably many exceptional null sets have a null
union \(N_0\); work outside \(N_0\). On \(\{S_c\leq n\}\), the time
\(S_c\) is finite and
\[
 A_{S_c\wedge n}=A_{S_c}=A_{T_c}+a_{S_c}\leq c+K
\]
by \cref{thm:p15-w02c-p15-021-exceedance-shift}. On \(\{S_c>n\}\),
the index \(n\) precedes the first exceedance, so
\(A_{S_c\wedge n}=A_n\leq c\). Hence \(A_{S_c\wedge n}\leq c+K\) almost
surely. By \cref{thm:p15-variation-stopping-transform} the stopped
martingale \(X_{S_c\wedge n}\) has predictable variation
\(A_{S_c\wedge n}\), and \cref{thm:p15-predictable-variation-identities}
applied to that stopped martingale gives
\[
 EX_{S_c\wedge n}^2=EA_{S_c\wedge n}\leq c+K .
\]
\end{proof}
\end{annexitem}

\begin{annexitem}{annex:p15-w02c-p15-021-unbounded-overshoot}{thm:p15-w02c-p15-021-unbounded-overshoot}{Computing both stops for the single-jump martingale}
\begin{proof}
\paragraph{The martingale and its variance.}
Write \(\Delta X_k=2^{k/2}\mathbf1_{\{J=k\}}\eta_k\) for \(k\geq1\). It
is \(\mathcal F_k\)-measurable, and
\[
 E|\Delta X_k|=2^{k/2}P(J=k)=2^{-k/2}<\infty,\qquad
 E(\Delta X_k)^2=2^kP(J=k)=1 .
\]
So each \(X_n\) is a finite sum of square-integrable variables and is
square integrable. The factor \(2^{k/2}\mathbf1_{\{J=k\}}\) is bounded and
\(\mathcal F_{k-1}\)-measurable, because \(J\) is
\(\mathcal F_0\)-measurable. The sign \(\eta_k\) is independent of
\(\mathcal F_{k-1}=\sigma(J,\eta_1,\ldots,\eta_{k-1})\), since the
variables \(J,\eta_1,\eta_2,\ldots\) are independent, and
\(E\eta_k=0\). Taking out what is known and using independence,
\[
 E(\Delta X_k\mid\mathcal F_{k-1})
   =2^{k/2}\mathbf1_{\{J=k\}}E\eta_k=0,
\]
so \(X\) is a martingale, and \(X_0=0\). Since \(\eta_k^2=1\),
\[
 E((\Delta X_k)^2\mid\mathcal F_{k-1})=2^k\mathbf1_{\{J=k\}} ,
\]
which is already finite, nonnegative and \(\mathcal F_{k-1}\)-measurable,
so it may be taken as \(a_k\). Summing, exactly one indicator
\(\mathbf1_{\{J=k\}}\) is nonzero, namely the one with \(k=J\), so
\(A_n=2^J\mathbf1_{\{J\leq n\}}\) and \(A_\infty=2^J\), which is finite
at every sample point. Finally
\[
 EA_n=\sum_{k=1}^n2^kP(J=k)=n,
\]
and \(EX_n^2=EA_n\) by \cref{thm:p15-predictable-variation-identities}.
In the same way \(X_n=2^{J/2}\eta_J\mathbf1_{\{J\leq n\}}\): the path is
zero before time \(J\) and constant from time \(J\) on.

\paragraph{The first-exceedance stop.}
Fix \(c>0\). Because \(A_n\in\{0,2^J\}\) and \(A_n=2^J\) exactly when
\(n\geq J\), one has \(S_c=J\) on \(\{2^J>c\}\) and \(S_c=\infty\) on
\(\{2^J\leq c\}\). On \(\{2^J\leq c\}\), \(S_c\wedge n=n\). On
\(\{2^J>c\}\), \(S_c\wedge n=J\wedge n\); if \(n\geq J\) then
\(X_{J}=X_n\) and \(A_J=A_n\) because both paths are constant from
time \(J\) on, and if \(n<J\) then \(J\wedge n=n\). In every case
\(X_{S_c\wedge n}=X_n\) and \(A_{S_c\wedge n}=A_n\). Therefore
\(EX_{S_c\wedge n}^2=EX_n^2=n\), which is unbounded in \(n\). This is
item~1.

On \(\{J\leq n,\ 2^J>c\}\), \(A_{S_c\wedge n}=A_n=2^J>c\). If
\(2^n>c\), that event contains \(\{J=n\}\), of probability
\(2^{-n}>0\). On \(\{S_c<\infty\}=\{2^J>c\}\), \(A_{S_c}=2^J\), so
\[
 E\bigl((A_{S_c}-c)\mathbf1_{\{S_c<\infty\}}\bigr)
   =\sum_{k:\,2^k>c}(2^k-c)2^{-k}
   =\sum_{k:\,2^k>c}\bigl(1-c2^{-k}\bigr).
\]
All but finitely many indices \(k\) satisfy \(2^k>2c\), and for those
the summand is at least \(1/2\); the series diverges. This is item~2.

\paragraph{The localizing stop.}
By \cref{thm:p15-w02c-p15-021-exceedance-shift}, \(T_c=S_c-1\), so
\(T_c=J-1\) on \(\{2^J>c\}\) and \(T_c=\infty\) on \(\{2^J\leq c\}\).
On \(\{2^J>c\}\), \(T_c\wedge n\leq J-1<J\), so
\(X_{T_c\wedge n}=0\) and \(A_{T_c\wedge n}=0\). On \(\{2^J\leq c\}\),
\(X_{T_c\wedge n}=X_n\) and \(A_{T_c\wedge n}=A_n\leq2^J\leq c\). Thus
\(X_{T_c\wedge n}=X_n\mathbf1_{\{2^J\leq c\}}\) and
\(A_{T_c\wedge n}\leq c\). The stopped variation identity of
\cref{thm:p15-variation-stopping-transform}, combined with
\cref{thm:p15-predictable-variation-identities}, gives
\(EX_{T_c\wedge n}^2=EA_{T_c\wedge n}\leq c\). This is item~3.
\end{proof}
\end{annexitem}
